\documentclass[11pt,a4paper]{article}
\usepackage{fontspec,amsmath,amssymb,amsthm,microtype}
\setmainfont{lmroman10-regular.otf}[BoldFont=lmroman10-bold.otf,ItalicFont=lmroman10-italic.otf,BoldItalicFont=lmroman10-bolditalic.otf]
\setmonofont{lmmono10-regular.otf}[Scale=MatchLowercase]
\usepackage[margin=25mm]{geometry}
\usepackage[unicode,hidelinks]{hyperref}
\setlength{\emergencystretch}{3em}
\setlength{\parskip}{0.4em}
\setlength{\parindent}{0pt}
\allowdisplaybreaks[2]
\raggedbottom
\title{\texorpdfstring{Continuous vector Brownian responses in L\ensuremath{{}^{2}}}{Continuous vector Brownian responses in L²}}
\author{LangevinLean mathematical exposition}
\date{}
\begin{document}\maketitle
\noindent Source module: \texttt{Kinetic\allowbreak{}Langevin/\allowbreak{}Vector\allowbreak{}Segment\allowbreak{}Equations.\allowbreak{}lean}. The accompanying HTML pairs this mathematical exposition with the exact Lean statements and records the scope of the formalization.


The scalar Brownian equations must also hold for Euclidean random vectors so they can drive the actual L\ensuremath{{}^{2}} kinetic diffusion used by the contraction proof. A continuous linear spatial assembly map achieves this directly: it preserves sums, scalar multiples, continuity, and Bochner time integrals.


\section{\texorpdfstring{Assemble scalar L\ensuremath{{}^{2}} variables through a continuous linear map}{Assemble scalar L² variables through a continuous linear map}}


For a scalar random variable \(Z\), spatial assembly produces the vector \((Z(\omega_i))_i\), evaluating it on independent spatial copies of the canonical source. The Lean definition is a finite sum. Each summand first pulls back along a measure-preserving coordinate evaluation, then inserts the scalar into one Euclidean basis coordinate. Both operations are continuous linear maps on L\ensuremath{{}^{2}}. The representative theorem is proved by the finite-sum almost-everywhere identity and the explicit representatives of these two maps; Euclidean extensionality verifies each coordinate. This construction avoids interchanging pointwise evaluation with an L\ensuremath{{}^{2}} Bochner integral.



{\raggedright\noindent Lean declarations: \texttt{spatial\allowbreak{}Lift}, \texttt{coe\_\allowbreak{}finset\_\allowbreak{}sum}, \texttt{spatial\allowbreak{}Lift\_\allowbreak{}coe}.\par}

\section{\texorpdfstring{Define vector responses and prove continuity}{Define vector responses and prove continuity}}


The vector response is the spatial assembly of the actual scalar response on \((a,a+t]\). The vector Brownian input is the assembly of the scalar Brownian increment. Their definitions remain on the existing canonical product source. The response representative theorem gives exactly the expected independent-coordinate vector. Continuity in time follows by composing the continuous scalar response with the continuous linear assembly map.



{\raggedright\noindent Lean declarations: \texttt{vector\allowbreak{}Segment\allowbreak{}Lp}, \texttt{vector\allowbreak{}Brownian\allowbreak{}Increment\allowbreak{}Lp}, \texttt{vector\allowbreak{}Segment\allowbreak{}Lp\_\allowbreak{}coe}, \texttt{continuous\_\allowbreak{}vector\allowbreak{}Segment\allowbreak{}Lp}.\par}

\section{\texorpdfstring{Lift the exact primitive and damped integral equations}{Lift the exact primitive and damped integral equations}}


The position response satisfies \(P^a_t=\int_0^tV^a_sds\), and the velocity response satisfies \(V^a_t=W^a_t-\gamma\int_0^tV^a_sds\). These equalities follow from the already proved scalar equations by applying the spatial assembly map. The proof commutes that continuous linear map with each Bochner integral using \texttt{interval\allowbreak{}Integral\_\allowbreak{}comp\_\allowbreak{}comm}. Thus the exact filtered Brownian input required by the generic mild-equation theorem is constructed for the vector L\ensuremath{{}^{2}} diffusion.



{\raggedright\noindent Lean declarations: \texttt{vector\allowbreak{}Segment\allowbreak{}Lp\_\allowbreak{}primitive}, \texttt{vector\allowbreak{}Segment\allowbreak{}Lp\_\allowbreak{}integral\_\allowbreak{}equation}.\par}

\section{\texorpdfstring{Identify the input with canonical vector Brownian increments}{Identify the input with canonical vector Brownian increments}}


The lifted input is proved almost everywhere equal to the vector whose coordinates are \(B_i(a+t)-B_i(a)\) for the existing canonical coordinate process. The scalar subtraction identity is pulled back through each measure-preserving spatial projection and combined with the spatial representative theorem. Continuity follows by identifying the increment with the unit-kernel vector response. This connects the L\ensuremath{{}^{2}} forcing to the actual canonical vector Brownian construction.



{\raggedright\noindent Lean declarations: \texttt{vector\allowbreak{}Brownian\allowbreak{}Increment\allowbreak{}Lp\_\allowbreak{}coe}, \texttt{continuous\_\allowbreak{}vector\allowbreak{}Brownian\allowbreak{}Increment\allowbreak{}Lp}.\par}

\section{\texorpdfstring{Identify the OU increment and numerical innovation endpoints}{Identify the OU increment and numerical innovation endpoints}}


The vector velocity response is exactly \(V^a_t=W^a_t-\gamma\int_0^te^{-\gamma(t-s)}W^a_sds\). Lean applies spatial assembly to the established scalar integration-by-parts formula, uses its linearity, and commutes it with the continuous Bochner integrand. This is precisely the OU input increment used by the kinetic splitting limit. At each grid interval, the paired L\ensuremath{{}^{2}} velocity and position response endpoints are proved almost everywhere equal to the pointwise \texttt{canonical\allowbreak{}KLMCInnovation} already equipped with the exact Gaussian law. The endpoint proof combines the two actual representative equalities with \texttt{Prod.\allowbreak{}ext}; it does not replace the constructed inputs by new independent Gaussian variables.



{\raggedright\noindent Lean declarations: \texttt{vector\allowbreak{}Segment\allowbreak{}Lp\_\allowbreak{}velocity\_\allowbreak{}eq\_\allowbreak{}brownian\_\allowbreak{}sub\_\allowbreak{}convolution}, \texttt{vector\allowbreak{}Segment\allowbreak{}Lp\_\allowbreak{}klmc\allowbreak{}Innovation\_\allowbreak{}coe}.\par}

\end{document}
